%============================================================
% CH2 · PATTERN 5 — EXERCISES, HINTS, ANSWERS
%============================================================
\subsection{Pattern 5: Exercises and Going Further}
\label{ssec:ch2-p5-ex}

\begin{patternbox}{Pattern 5 Summary Card}
\renewcommand{\arraystretch}{1.25}
\begin{tabularx}{\linewidth}{@{}>{\bfseries}l X@{}}
Recognise & many starts, few targets: ``from which cells / vertices can you reach \dots?'', ``every good one must escape'', ``touches the border'', ``ancestors'', ``forever'' \\
Reverse & directed: store $a \to b$ in \texttt{radj[b]}; directed rule in a grid: flip the comparison; undirected: only the start changes \\
Many targets & push all of them before the loop --- one search, \VE \\
Per vertex & one reverse search per vertex, \bigO{n(n+m)}, sorted lists by an index-order scan \\
Good vs.\ bad & neutralise the bad starts with a change every solution needs, then one search from the target \\
Forever & peel out-degree-0 vertices along reversed edges (Kahn on $G^R$) \\
Top traps & \texttt{radj[a]} instead of \texttt{radj[b]}; half-reversed grid rule; targets counted as answers; a walled exit with nobody to save \\
\end{tabularx}
\end{patternbox}

\exlegend
\begin{enumerate}[leftmargin=*, label=\textbf{E5.\arabic*}, itemsep=6pt]
    \item \exeasy\ \pclassic \lc{Medium}\ \exprob{Surrounded Regions}{LeetCode 130}{https://leetcode.com/problems/surrounded-regions/}. A board of \texttt{X} and \texttt{O}. Every region of \texttt{O} cells (4-connected) that does not touch the border is captured: flip it to \texttt{X}. Modify the board in place. ($1 \le R, C \le 200$.)
    \item \exmedium\ \pclassic \lc{Medium}\ \exprob{Pacific Atlantic Water Flow}{LeetCode 417}{https://leetcode.com/problems/pacific-atlantic-water-flow/}. An $R \times C$ height map; the Pacific touches the top and left edges, the Atlantic the bottom and right edges. Rain flows to a side-neighbour of height $\le$ the current one, and from a border cell into the adjacent ocean. List all cells from which water can reach both oceans. ($R, C \le 200$.)
    \item \exmedium\ \ptricky \cf{1600}\ \exprob{Lakes in Berland}{Codeforces 723D}{https://codeforces.com/problemset/problem/723/D}. A map of water (\texttt{.}) and land (\texttt{*}) surrounded by the ocean. A lake is a maximal 4-connected water region with no cell on the border. Fill as few water cells as possible so that exactly $k$ lakes remain; print the number and the new map. ($n, m \le 50$.)
    \item \exmedium\ \pconcept \exprob{Endless Walk}{AtCoder ABC245 F}{https://atcoder.jp/contests/abc245/tasks/abc245_f}. A simple directed graph with $N$ vertices and $M$ edges. Count the vertices from which one can walk along edges forever (choosing the edges cleverly). ($N, M \le 2 \cdot 10^5$.)
    \item \exmedium\ \pconcept \exown{Who can reach $T$?} Given a directed graph and a set $T$ of target vertices, mark every vertex that can reach some vertex of $T$. Implement the multi-target reverse search and the naive ``one forward search per vertex'', check that they agree on random graphs, and measure both for $n = 10^3, 10^4, 10^5$ with $m = 2n$ random edges and $|T| = 3$.
    \item \exhard\ \precent \lc{Medium}\ \exprob{Minimize the Maximum Edge Weight of Graph}{LeetCode 3419}{https://leetcode.com/problems/minimize-the-maximum-edge-weight-of-graph/}. A directed weighted graph on $n$ vertices and an integer \texttt{threshold} $\ge 1$. Remove edges so that vertex 0 is reachable from every vertex and every vertex keeps at most \texttt{threshold} outgoing edges, minimising the largest remaining weight. Return it, or $-1$. ($n, m \le 10^5$.)
\end{enumerate}

\subsubsection*{Hints}
\begin{enumerate}[leftmargin=*, label=\textbf{E5.\arabic*}]
    \item Which \texttt{O} cells are \emph{safe}? Find those first, starting from the border.
    \item Asking ``where does water from this cell go?'' for every cell is $R \cdot C$ searches. Ask instead ``which cells can send water to the Pacific?'' --- once, from the coast, uphill.
    \item The water regions that touch the border are not lakes; find them from the border (or flag them during a flood fill). Which lakes are cheapest to fill?
    \item A vertex with no outgoing edge is hopeless. What happens to a vertex all of whose edges lead to hopeless vertices?
    \item Template 1 of the pattern. For the naive version, stop each search as soon as it finds a target --- it is still slow.
    \item Binary search on the answer $W$. With only the edges of weight $\le W$: does everybody reach 0? And is \texttt{threshold} ever the real obstacle?
\end{enumerate}

\subsubsection*{Answers}

\begin{solutionbox}[E5.1 Surrounded Regions]
\textbf{Reverse the roles.} ``Is this region captured?'' per region becomes one search from all \texttt{O} cells on the border: everything it reaches is safe (temporarily marked \texttt{S}); every \texttt{O} left unmarked is captured. Finally \texttt{S} $\to$ \texttt{O}, \texttt{O} $\to$ \texttt{X}.
\begin{lstlisting}[style=cpp]
#include <bits/stdc++.h>
using namespace std;

class Solution {
public:
    void solve(vector<vector<char>>& board) {
        int R = board.size(), C = board[0].size();
        int dr[] = {1, -1, 0, 0}, dc[] = {0, 0, 1, -1};
        vector<pair<int, int>> st;
        for (int r = 0; r < R; r++)
            for (int c = 0; c < C; c++) {
                bool border = (r == 0 || c == 0 || r == R - 1 || c == C - 1);
                if (border && board[r][c] == 'O') {
                    board[r][c] = 'S'; st.push_back({r, c});
                }
            }
        while (!st.empty()) {                     // S = safe: connected to the border
            int r = st.back().first, c = st.back().second; st.pop_back();
            for (int d = 0; d < 4; d++) {
                int nr = r + dr[d], nc = c + dc[d];
                if (nr < 0 || nc < 0 || nr >= R || nc >= C) continue;
                if (board[nr][nc] != 'O') continue;
                board[nr][nc] = 'S'; st.push_back({nr, nc});
            }
        }
        for (auto& row : board)
            for (char& ch : row) ch = (ch == 'S') ? 'O' : 'X';   // captured O -> X
    }
};
\end{lstlisting}
Example 1 \texttt{XXXX / XOOX / XXOX / XOXX} $\to$ \texttt{XXXX / XXXX / XXXX / XOXX} (the bottom \texttt{O} touches the border and survives); example 2 \texttt{X} $\to$ \texttt{X}. Checked against a brute force (one search per \texttt{O} cell) on 5000 random boards up to $6 \times 6$. \bigO{RC}; iterative, since a $200 \times 200$ snake of \texttt{O} would make a recursive DFS $\approx 20\,000$ levels deep.
\end{solutionbox}

\begin{solutionbox}[E5.2 Pacific Atlantic Water Flow]
Water flows from $a$ to a neighbour $b$ if $h_b \le h_a$. Reversed: from $b$ we may \textbf{climb} to $a$ if $h_a \ge h_b$. Run this climb from all Pacific border cells and, separately, from all Atlantic border cells; the answer is the intersection.
\begin{center}
\begin{tikzpicture}[x=6.5mm, y=6.5mm]
    \foreach \r/\row in {0/{1,2,2,3,5}, 1/{3,2,3,4,4}, 2/{2,4,5,3,1}, 3/{6,7,1,4,5}, 4/{5,1,1,2,4}}
        \foreach \h [count=\c from 0] in \row \node[g2p5-e] at (\c,-\r) {\h};
    \foreach \r/\c in {0/4, 1/3, 1/4, 2/2, 3/0, 3/1, 4/0}
        \node[g2p5-e, fill=answerfill, draw=accentorange, line width=1.1pt] at (\c,-\r) {};
    \foreach \r/\row in {0/{1,2,2,3,5}, 1/{3,2,3,4,4}, 2/{2,4,5,3,1}, 3/{6,7,1,4,5}, 4/{5,1,1,2,4}}
        \foreach \h [count=\c from 0] in \row \node[font=\scriptsize\bfseries] at (\c,-\r) {\h};
    \draw[primaryblue, line width=2pt] (-0.55,-4.5) -- (-0.55,0.55) -- (4.5,0.55) node[right, font=\scriptsize] {Pacific};
    \draw[successgreen, line width=2pt] (-0.5,-4.55) -- (4.55,-4.55) -- (4.55,0.5);
    \node[font=\scriptsize, anchor=west] at (4.7,-4.4) {Atlantic};
\end{tikzpicture}
\end{center}
\begin{lstlisting}[style=cpp]
#include <bits/stdc++.h>
using namespace std;

class Solution {
    int R, C;
    // reverse search from an ocean's border: climb to neighbours that are NOT LOWER
    vector<vector<char>> climb(const vector<vector<int>>& h, vector<pair<int, int>> st) {
        vector<vector<char>> ok(R, vector<char>(C, 0));
        for (auto& p : st) ok[p.first][p.second] = 1;
        int dr[] = {1, -1, 0, 0}, dc[] = {0, 0, 1, -1};
        while (!st.empty()) {
            int r = st.back().first, c = st.back().second; st.pop_back();
            for (int d = 0; d < 4; d++) {
                int nr = r + dr[d], nc = c + dc[d];
                if (nr < 0 || nc < 0 || nr >= R || nc >= C || ok[nr][nc]) continue;
                if (h[nr][nc] < h[r][c]) continue;     // no flow from nr up to r
                ok[nr][nc] = 1; st.push_back({nr, nc});
            }
        }
        return ok;
    }
public:
    vector<vector<int>> pacificAtlantic(vector<vector<int>>& heights) {
        R = heights.size(); C = heights[0].size();
        vector<pair<int, int>> pac, atl;
        for (int r = 0; r < R; r++) { pac.push_back({r, 0}); atl.push_back({r, C - 1}); }
        for (int c = 0; c < C; c++) { pac.push_back({0, c}); atl.push_back({R - 1, c}); }
        sort(pac.begin(), pac.end());              // corners appear twice: deduplicate
        pac.erase(unique(pac.begin(), pac.end()), pac.end());
        sort(atl.begin(), atl.end());
        atl.erase(unique(atl.begin(), atl.end()), atl.end());
        vector<vector<char>> P = climb(heights, pac), A = climb(heights, atl);
        vector<vector<int>> res;
        for (int r = 0; r < R; r++)
            for (int c = 0; c < C; c++)
                if (P[r][c] && A[r][c]) res.push_back({r, c});
        return res;
    }
};
\end{lstlisting}
Example 1 gives \texttt{[0,4] [1,3] [1,4] [2,2] [3,0] [3,1] [4,0]} (orange above), example 2 \texttt{[[1]]} gives \texttt{[0,0]}; both equal the official outputs. A brute force that floods downhill from every cell agreed on 5000 random maps up to $6 \times 6$ with heights $0..4$ (many ties, which is where the $\le$ vs.\ $<$ bug lives). \bigO{RC}.
\end{solutionbox}

\begin{solutionbox}[E5.3 Lakes in Berland]
\textbf{The trap:} water connected to the ocean is not a lake, even if it is large, and filling it is never needed. Flood-fill every water component and flag those with a border cell (equivalently: first search from all border water cells, then the unmarked components are the lakes). With $L$ lakes, fill the $L - k$ \emph{smallest} ones: each lake must be filled completely to disappear (filling part of it leaves a smaller lake, or splits it into more lakes), so the cost is the sum of the sizes of the lakes we remove, minimised by removing the smallest.
\begin{lstlisting}[style=cpp]
#include <bits/stdc++.h>
using namespace std;

int n, m, k;
vector<string> g;
const int dr[] = {1, -1, 0, 0}, dc[] = {0, 0, 1, -1};

// iterative flood fill of '.' cells from (r, c); border = some cell is on the edge
vector<pair<int, int>> fill(int r, int c, vector<vector<char>>& seen, bool& border) {
    vector<pair<int, int>> cells, st(1, make_pair(r, c));
    seen[r][c] = 1; border = false;
    while (!st.empty()) {
        int x = st.back().first, y = st.back().second; st.pop_back();
        cells.push_back(make_pair(x, y));
        if (x == 0 || y == 0 || x == n - 1 || y == m - 1) border = true;
        for (int d = 0; d < 4; d++) {
            int a = x + dr[d], b = y + dc[d];
            if (a < 0 || b < 0 || a >= n || b >= m) continue;
            if (seen[a][b] || g[a][b] != '.') continue;
            seen[a][b] = 1; st.push_back(make_pair(a, b));
        }
    }
    return cells;
}

int main() {
    cin >> n >> m >> k;
    g.resize(n);
    for (auto& s : g) cin >> s;
    vector<vector<char>> seen(n, vector<char>(m, 0));
    vector<vector<pair<int, int>>> lakes;
    for (int r = 0; r < n; r++)
        for (int c = 0; c < m; c++)
            if (g[r][c] == '.' && !seen[r][c]) {
                bool border;
                vector<pair<int, int>> cells = fill(r, c, seen, border);
                if (!border) lakes.push_back(cells);   // touches the ocean: not a lake
            }
    sort(lakes.begin(), lakes.end(),
         [](const vector<pair<int, int>>& a, const vector<pair<int, int>>& b) {
             return a.size() < b.size();
         });
    int filled = 0;
    for (size_t i = 0; i + k < lakes.size(); i++)      // fill all but the k largest
        for (auto& p : lakes[i]) { g[p.first][p.second] = '*'; filled++; }
    cout << filled << "\n";
    for (auto& s : g) cout << s << "\n";
}
\end{lstlisting}
Both samples are reproduced exactly (\texttt{1} and the maps). In sample 1 the bottom-left water touches the border and is ignored, although it is as large as the lake of size 2. A brute force that tries all sets of water cells to fill (grids up to $5 \times 5$ with $\le 12$ water cells, random $k$) agreed on 1500 tests: same minimum, a map with exactly $k$ lakes, only water cells changed.
\end{solutionbox}

\begin{solutionbox}[E5.4 Endless Walk]
A vertex with out-degree 0 cannot start an endless walk; delete it, and decrement the out-degree of its predecessors --- found through the reversed edges. Repeat. The survivors all have an edge to another survivor, so a walk can go on forever from each of them; every deleted vertex leads only into deleted vertices. This is template 4 of the pattern (Kahn's algorithm on $G^R$).
\begin{center}
\begin{tikzpicture}[x=1.2cm]
    \node[gvsel] (w1) at (0,0) {1}; \node[gvsel] (w2) at (1.3,0) {2}; \node[gvsel] (w3) at (2.6,0.6) {3};
    \node[gvsel] (w4) at (2.6,-0.6) {4}; \node[gvdead] (w5) at (3.9,-0.6) {5};
    \draw[gde] (w1) -- (w2); \draw[gde] (w2) -- (w3); \draw[gde] (w3) -- (w4); \draw[gde] (w4) -- (w2);
    \draw[gfaded, -{Stealth[length=2.2mm]}] (w4) -- (w5);
    \node[note, text width=5.6cm, align=left, anchor=west] at (4.6,0) {Sample 1: only 5 has out-degree 0. Deleting it leaves 4 with out-degree 1 (edge $4 \to 2$), so the peeling stops: answer \textbf{4}.};
\end{tikzpicture}
\end{center}
\begin{lstlisting}[style=cpp]
#include <bits/stdc++.h>
using namespace std;

int main() {
    ios::sync_with_stdio(false);
    cin.tie(nullptr);
    int n, m;
    cin >> n >> m;
    vector<vector<int>> radj(n + 1);
    vector<int> out(n + 1, 0);
    for (int i = 0; i < m; i++) {
        int u, v;
        cin >> u >> v;
        radj[v].push_back(u);              // remember who points INTO v
        out[u]++;
    }
    vector<int> q;
    for (int v = 1; v <= n; v++) if (out[v] == 0) q.push_back(v);
    int removed = 0;
    while (!q.empty()) {
        int v = q.back(); q.pop_back();
        removed++;                         // every walk from v is finite
        for (int u : radj[v])
            if (--out[u] == 0) q.push_back(u);   // u's last exit led to a dead end
    }
    cout << n - removed << "\n";
}
\end{lstlisting}
Samples: \texttt{4} and \texttt{2} (sample 2: $1 \rightleftarrows 2$ and an isolated 3). A brute force with transitive closure (a vertex counts iff it reaches a vertex that lies on a cycle) agreed on 2000 random graphs with $N \le 8$. \bigO{N + M}, no recursion. A DFS alternative: colour the graph and count the vertices that reach a back edge --- correct, but it needs the iterative DFS and a second pass; the peeling is simpler.
\end{solutionbox}

\begin{solutionbox}[E5.5 Who can reach $T$?]
\begin{lstlisting}[style=cpp]
#include <bits/stdc++.h>
using namespace std;

// (a) one multi-target search in G^R
vector<char> reverseSearch(const vector<vector<int>>& radj, const vector<int>& T) {
    vector<char> can(radj.size(), 0);
    vector<int> st;
    for (int t : T) if (!can[t]) { can[t] = 1; st.push_back(t); }
    while (!st.empty()) {
        int v = st.back(); st.pop_back();
        for (int u : radj[v]) if (!can[u]) { can[u] = 1; st.push_back(u); }
    }
    return can;
}

// (b) a forward search from every vertex, stopped at the first target
vector<char> forwardEach(const vector<vector<int>>& adj, const vector<char>& inT,
                         long long& work) {
    int n = adj.size();
    vector<char> can(n, 0);
    vector<int> mark(n, -1), st;
    for (int s = 0; s < n; s++) {
        mark[s] = s; st.assign(1, s);
        while (!st.empty() && !can[s]) {
            int u = st.back(); st.pop_back(); work++;
            if (inT[u]) can[s] = 1;
            for (int v : adj[u]) if (mark[v] != s) { mark[v] = s; st.push_back(v); }
        }
    }
    return can;
}

int main() {
    mt19937 rng(11);
    int checks = 0, bad = 0;
    for (int t = 0; t < 2000; t++) {                 // small graphs: both must agree
        int n = rng() % 9 + 1, m = rng() % 20;
        vector<vector<int>> adj(n), radj(n);
        for (int i = 0; i < m; i++) {
            int a = rng() % n, b = rng() % n;
            adj[a].push_back(b); radj[b].push_back(a);
        }
        vector<int> T; vector<char> inT(n, 0);
        for (int v = 0; v < n; v++) if (rng() % 4 == 0) { T.push_back(v); inT[v] = 1; }
        long long w = 0;
        bad += reverseSearch(radj, T) != forwardEach(adj, inT, w);
        checks++;
    }
    cout << checks << " small graphs, " << bad << " differences\n";
    for (int n : {1000, 10000, 100000}) {             // timing: m = 2n, |T| = 3
        vector<vector<int>> adj(n), radj(n);
        for (int i = 0; i < 2 * n; i++) {
            int a = rng() % n, b = rng() % n;
            adj[a].push_back(b); radj[b].push_back(a);
        }
        vector<int> T = {0, 1, 2};
        vector<char> inT(n, 0);
        inT[0] = inT[1] = inT[2] = 1;
        long long work = 0;
        clock_t t0 = clock();
        vector<char> a = reverseSearch(radj, T);
        clock_t t1 = clock();
        vector<char> b = forwardEach(adj, inT, work);
        clock_t t2 = clock();
        cout << "n=" << n << ": reach T " << count(a.begin(), a.end(), 1)
             << ", equal " << (a == b)
             << ", reverse " << (t1 - t0) * 1000 / CLOCKS_PER_SEC << " ms, forward "
             << (t2 - t1) * 1000 / CLOCKS_PER_SEC << " ms (" << work << " pops)\n";
    }
}
\end{lstlisting}
Output on our machine:
\begin{lstlisting}[style=trace]
2000 small graphs, 0 differences
n=1000: reach T 805, equal 1, reverse 0 ms, forward 1 ms (59355 pops)
n=10000: reach T 7993, equal 1, reverse 0 ms, forward 724 ms (14748395 pops)
n=100000: reach T 79599, equal 1, reverse 8 ms, forward 63193 ms (384723015 pops)
\end{lstlisting}
About 80\% of the vertices reach $T$ (with $m = 2n$ most vertices lie in or before the giant strongly connected part). The forward searches that \emph{fail} --- the remaining 20\% --- are the expensive ones: each explores everything its start can reach before giving up, and that is a large part of the graph. Early stopping does not save the naive method; the reverse search is about $8000$ times faster at $n = 10^5$ (earlier runs of the same experiment: 10--12\,ms against 64--95\,s) and its cost grows linearly.
\end{solutionbox}

\begin{solutionbox}[E5.6 Minimize the Maximum Edge Weight of Graph]
\textbf{Two observations.} (1) ``0 is reachable from every vertex'' is one search from 0 in $G^R$. (2) The \texttt{threshold} is a decoy: if everybody reaches 0, the reverse search tree gives every vertex $\ne 0$ exactly one outgoing edge (towards its parent, i.e.\ towards 0) and vertex 0 none; keep only those $n - 1$ edges, and every out-degree is $\le 1 \le$ \texttt{threshold}. Removing edges never hurts the maximum weight. So the answer is the smallest $W$ for which the edges of weight $\le W$ let everybody reach 0 --- monotone in $W$, hence a binary search over the distinct weights, each step one reverse search.
\begin{center}
\begin{tikzpicture}[x=1.25cm]
    \node[gvdst] (z0) at (0,0) {0}; \node[gv] (z4) at (1.4,0) {4}; \node[gv] (z3) at (2.8,0) {3};
    \node[gv] (z2) at (2.8,-1.2) {2}; \node[gv] (z1) at (1.4,-1.2) {1};
    \draw[g2p5-rt] (z0) -- node[gw, above]{1} (z4);
    \draw[g2p5-rt] (z4) -- node[gw, above]{2} (z3);
    \draw[g2p5-rt] (z3) -- node[gw, right]{2} (z2);
    \draw[g2p5-rt] (z2) -- node[gw, below]{1} (z1);
    \draw[gfaded, -{Stealth[length=2.2mm]}] (z1) -- node[gw, left, pos=0.4]{3} (z3);
    \draw[gfaded, -{Stealth[length=2.2mm]}] (z1) -- node[gw, left]{5} (z4);
    \node[note, text width=5.0cm, align=left, anchor=west] at (3.6,-0.6) {Example 3 with $W = 2$: the reverse search from 0 reaches $4, 3, 2, 1$ (blue); the edges $1 \to 3$ (3) and $1 \to 4$ (5) are dropped and every vertex keeps one outgoing edge. With $W = 1$ vertex 3 is cut off. Answer \textbf{2}.};
\end{tikzpicture}
\end{center}
\begin{lstlisting}[style=cpp]
#include <bits/stdc++.h>
using namespace std;

class Solution {
    // only edges of weight <= W: does every vertex reach 0?  = one search from 0 in G^R
    bool allReachZero(int n, const vector<vector<pair<int, int>>>& radj, int W) {
        vector<char> seen(n, 0);
        vector<int> st(1, 0);
        seen[0] = 1;
        int cnt = 1;
        while (!st.empty()) {
            int v = st.back(); st.pop_back();
            for (auto& e : radj[v])                   // original edge e.first -> v
                if (e.second <= W && !seen[e.first]) {
                    seen[e.first] = 1; cnt++; st.push_back(e.first);
                }
        }
        return cnt == n;
    }
public:
    int minMaxWeight(int n, vector<vector<int>>& edges, int threshold) {
        vector<vector<pair<int, int>>> radj(n);
        vector<int> ws;
        for (auto& e : edges) {
            radj[e[1]].push_back({e[0], e[2]});
            ws.push_back(e[2]);
        }
        sort(ws.begin(), ws.end());
        ws.erase(unique(ws.begin(), ws.end()), ws.end());
        if (!allReachZero(n, radj, ws.back())) return -1;
        int lo = 0, hi = ws.size() - 1;               // smallest index that works
        while (lo < hi) {
            int mid = (lo + hi) / 2;
            if (allReachZero(n, radj, ws[mid])) hi = mid; else lo = mid + 1;
        }
        return ws[lo];                                // threshold never matters
    }
};
\end{lstlisting}
The four examples give \texttt{1 -1 2 -1}, as required. A brute force over all subsets of at most 10 edges (checking out-degrees and reachability directly) agreed on 3000 random graphs with $n \le 6$ and random thresholds. \bigO{(n+m)\log m}, iterative (a path of $10^5$ vertices is a legal input). Alternative without binary search: a Prim/Dijkstra-like ``minimax'' search from 0 in $G^R$ that always takes the cheapest edge leaving the reached set; the answer is the largest edge it takes.
\end{solutionbox}
